% !TeX root = ../main.tex
\chapter{Operational and Integration Details}
\label{appendix:operational-details}

This appendix collects technical details that are useful for
implementation review but too specific for the main Chapter~4 design
narrative.

\section{gRPC Internal API Details}
\label{appendix:grpc-proto-details}

\textbf{Protocol boundary.} The system uses three distinct communication
channels, each with a clear purpose:

\begin{itemize}
    \item \textbf{HTTP (REST + SSE):} only between the browser and
          \texttt{web-bff}, because browsers cannot speak gRPC directly.
          SSE is also used to stream chat tokens and pipeline progress
          back to the browser.
    \item \textbf{gRPC:} all synchronous backend-to-backend service
          calls, leveraging Protocol Buffers for cross-language type
          safety (TypeScript NestJS, Python FastAPI).
    \item \textbf{BullMQ events on Redis:} all asynchronous worker
          interactions; workers never call each other directly.
\end{itemize}

\textbf{Diagram convention.} Sequence and dataflow diagrams in
Chapter~\ref{chap:phan-tich-thiet-ke} use HTTP-style shorthand (e.g.
\texttt{POST /documents/upload-session}) for readability, but the
underlying call between \texttt{web-bff} and \texttt{core-api} is a
gRPC method invocation
(\texttt{CoreApiService.CreateUploadSession}). The mapping between the
two notations is one-to-one: each REST verb in a diagram corresponds
to exactly one gRPC method on the service contract.

\textbf{Service contracts.} Two gRPC services are defined:

\begin{itemize}
    \item \textbf{\texttt{CoreApiService}} (served by \texttt{core-api}):
          the main business surface, called by \texttt{web-bff} for
          resource CRUD (\texttt{CreateUploadSession},
          \texttt{ConfirmUpload}, \texttt{ListDocuments},
          \texttt{PublishContent}) and by \texttt{chat-service} for
          permission/context resolution (\texttt{GetUserContext},
          \texttt{CheckCourseAccess}, \texttt{ListLOByCourse},
          \texttt{GetChunkMetadata}). The latter is called once at SSE
          handshake to populate the JWT ticket with
          \texttt{allowed\_lo\_ids} and \texttt{allowed\_document\_ids}
          so subsequent retrieval is scoped without further round trips.
    \item \textbf{\texttt{ChatService}} (served by \texttt{chat-service}):
          synchronous query surface called by \texttt{core-api}
          (\texttt{GetChatHistory}, \texttt{ListChatSessions},
          \texttt{SearchVideos}) for analytics dashboards and audit.
          The chat token stream itself is \emph{not} a gRPC method; it is
          exposed on \texttt{/chat/stream} as Server-Sent Events because
          the browser consumes it directly through \texttt{web-bff}.
\end{itemize}

\begin{lstlisting}[caption=Representative internal gRPC proto contract]
syntax = "proto3";
package ai_lms.v1;

service CoreApiService {
  rpc CreateUploadSession(CreateUploadSessionRequest)
      returns (CreateUploadSessionResponse);
  rpc ConfirmUpload(ConfirmUploadRequest)
      returns (DocumentSummary);
  rpc ListDocuments(ListDocumentsRequest)
      returns (ListDocumentsResponse);
  rpc PublishContent(PublishContentRequest)
      returns (PublishContentResponse);
  rpc GetUserContext(GetUserContextRequest)
      returns (UserContext);
  rpc CheckCourseAccess(CheckCourseAccessRequest)
      returns (AccessDecision);
  rpc ListLOByCourse(ListLOByCourseRequest)
      returns (ListLOByCourseResponse);
  rpc GetChunkMetadata(GetChunkMetadataRequest)
      returns (ChunkMetadata);
}

service ChatService {
  rpc GetChatHistory(GetChatHistoryRequest)
      returns (GetChatHistoryResponse);
  rpc ListChatSessions(ListChatSessionsRequest)
      returns (ListChatSessionsResponse);
  rpc SearchVideos(SearchVideosRequest)
      returns (SearchVideosResponse);
}

message CreateUploadSessionRequest {
  string course_id = 1;
  string file_name = 2;
  string mime_type = 3;
  int64 size_bytes = 4;
}

message CreateUploadSessionResponse {
  string document_id = 1;
  string upload_url = 2;
  string storage_key = 3;
  PipelineStatus status = 4;
}

message UserContext {
  string user_id = 1;
  string course_id = 2;
  repeated Role roles = 3;
  repeated string allowed_document_ids = 4;
  repeated string allowed_lo_ids = 5;
}

message Citation {
  string source_id = 1;
  string source_type = 2;
  string title = 3;
  int32 page_number = 4;
  int32 start_ms = 5;
  int32 end_ms = 6;
}

enum Role {
  ROLE_UNSPECIFIED = 0;
  ROLE_LEARNER = 1;
  ROLE_INSTRUCTOR = 2;
  ROLE_ADMIN = 3;
}

enum PipelineStatus {
  PIPELINE_STATUS_UNSPECIFIED = 0;
  QUEUED = 1;
  PARSING = 2;
  CHUNKING = 3;
  EMBEDDING = 4;
  INDEXED = 5;
  ERROR = 6;
}
\end{lstlisting}

\textbf{Shared types.} \texttt{shared.proto} defines common messages
(\texttt{Citation}, \texttt{Role} enum, \texttt{PipelineStatus} enum)
imported by both services, ensuring the same vocabulary across
runtimes.

\textbf{Authentication.} Each gRPC request carries the internal JWT in
the \texttt{authorization} metadata header. The server verifies the
token signature using the BFF's public key (for web-bff $\to$ core-api)
or core-api's public key (for chat-service $\to$ core-api). Token
scopes determine whether the call is permitted.

\textbf{What gRPC does not carry.} Asynchronous worker results are
never delivered via gRPC. Workers write to PostgreSQL (source of truth)
and publish BullMQ events on Redis; consumers learn about new state by
subscribing to events or polling the database. Pipeline progress
updates also avoid gRPC: workers publish through Redis Pub/Sub,
\texttt{web-bff} subscribes, and the browser receives updates via SSE.
This separation removes temporal coupling between long-running jobs and
the request/response services.

\section{Observability Details}
\label{appendix:observability-details}

\subsection{Structured Logging}

Each runtime emits JSON logs through \texttt{structlog} (Python) or
\texttt{pino} (NestJS). Required fields:

\begin{itemize}
    \item \texttt{timestamp}: ISO8601.
    \item \texttt{level}: \texttt{debug}/\texttt{info}/\texttt{warn}/\texttt{error}.
    \item \texttt{service}: runtime name.
    \item \texttt{trace\_id}, \texttt{span\_id}: OpenTelemetry standard,
          automatically injected into each log record by SDK.
    \item \texttt{user\_id}, \texttt{course\_id}: if business context
          exists.
    \item \texttt{message}: human-readable.
    \item \texttt{extra}: structured payload (input size, duration,
          error\_class).
\end{itemize}

Logs are aggregated by Loki (lightweight) or Elasticsearch (production).
Grafana provides query UI and allows jumping directly from a log entry to
the corresponding distributed trace through the shared \texttt{trace\_id}.

\textbf{Log sampling:} \texttt{debug} level is enabled only for sampled
traces (\texttt{traceflags=01}) or when the client attaches baggage
\texttt{debug=1}, preventing log volume explosion.

\subsection{Metrics Collection}

Prometheus scrapes each service's \texttt{/metrics} endpoint every 15s.
Metrics are grouped into three categories.

\textbf{Application metrics:}
\begin{itemize}
    \item \texttt{http\_request\_duration\_seconds} (histogram, labels
          method/route/status).
    \item \texttt{http\_requests\_total} (counter).
    \item \texttt{grpc\_request\_duration\_seconds}.
    \item \texttt{job\_duration\_seconds} (per worker queue).
    \item \texttt{job\_queue\_depth} (gauge).
    \item \texttt{llm\_call\_duration\_seconds}, \texttt{llm\_token\_total}.
    \item \texttt{embedding\_duration\_seconds}.
    \item \texttt{vector\_search\_duration\_seconds}.
\end{itemize}

\textbf{Business metrics:}
\begin{itemize}
    \item \texttt{document\_uploaded\_total}, \texttt{video\_uploaded\_total}.
    \item \texttt{content\_generated\_total} (label type=card/quiz).
    \item \texttt{content\_published\_total},
          \texttt{content\_rejected\_total}.
    \item \texttt{quiz\_attempts\_total}, \texttt{chat\_messages\_total}.
    \item \texttt{lo\_coverage\_ratio} (gauge per course).
\end{itemize}

\textbf{System metrics:} CPU, RAM, disk, and network of containers/host
(node\_exporter, cAdvisor).

Grafana dashboards contain panels for each service plus golden signals
(latency, throughput, error rate, saturation).

\subsection{Distributed Tracing}

OpenTelemetry SDK is used in each service, with automatic instrumentation
for HTTP/gRPC/SQL clients. Traces propagate through the
\texttt{traceparent} header.

\textbf{Main spans:}
\begin{itemize}
    \item HTTP request entry span.
    \item Use case execution span.
    \item Adapter call span (DB, Qdrant, Neo4j, LLM, MinIO).
    \item Worker job span.
\end{itemize}

Traces are exported through OTLP to Tempo (lightweight) or Jaeger. Trace
ID is attached to logs for cross-reference during debugging.

\subsection{AI Cost Monitoring}

LLM calls are the largest cost source. The system tracks:

\begin{itemize}
    \item \texttt{llm\_token\_input\_total}, \texttt{llm\_token\_output\_total}
          (counter, labels model/provider/use\_case).
    \item \texttt{llm\_cost\_usd\_total} (counter, computed from provider
          pricing table).
    \item Per-user, per-course, per-day breakdown in PostgreSQL table
          \texttt{llm\_usage\_logs}.
\end{itemize}

Cost spike alert: Prometheus rule
\texttt{rate(llm\_cost\_usd\_total[1h]) > threshold}.

\textbf{Cost optimization:}
\begin{itemize}
    \item Cache prompt results when input hash is identical for
          deterministic generation.
    \item Use smaller models (Gemini Flash) for tasks that do not require
          large models.
    \item Limit context window: send only top-N chunks rather than the
          full document.
\end{itemize}

\subsection{Error Tracking}

Sentry integrates with \texttt{web-bff} and \texttt{core-api}:

\begin{itemize}
    \item Capture exceptions with stack trace + context (user, route).
    \item Group by fingerprint and deduplicate.
    \item Alert by email/Slack when error rate spikes.
\end{itemize}

Worker exceptions are structured-logged and pushed into the dead letter
queue, with an internal endpoint for inspection and requeue.

\section{Deployment Details}
\label{appendix:deployment-details}

\subsection{Docker Compose Topology}

The project is deployed with Docker Compose using two stacks:

\textbf{Stack \texttt{ai-system}} (main system):
\begin{itemize}
    \item Services: \texttt{web-bff}, \texttt{core-api},
          \texttt{chat-service}, \texttt{document-worker},
          \texttt{enrichment-worker}, \texttt{video-worker}.
    \item Data: \texttt{postgres}, \texttt{redis}, \texttt{qdrant},
          \texttt{neo4j}, \texttt{minio}.
    \item Observability: \texttt{prometheus}, \texttt{grafana},
          \texttt{tempo}, \texttt{loki}.
    \item Total: around 15 containers.
\end{itemize}

\textbf{Stack \texttt{lms}} (optional in demo/dev environment):
\begin{itemize}
    \item Canvas LMS is used as the LTI Platform in internal demo
          scenarios; it can be Canvas Cloud or self-hosted Canvas Open
          Source.
    \item If self-hosted Canvas is used, LMS has its own database and does
          not share a database with \texttt{ai-system}.
\end{itemize}

In production or real pilot environments, Canvas is usually an external
system; only \texttt{web-bff} needs to expose a public HTTPS endpoint to
receive OIDC/LTI launches from Canvas. In self-hosted demos, the two
stacks can communicate through the shared Docker network
\texttt{lms\_shared}.

\subsection{Network and Volume}

\textbf{Network:}
\begin{itemize}
    \item \texttt{ai\_internal}: internal network among AI services, not
          exposed externally.
    \item \texttt{lms\_shared}: bridge between \texttt{ai-system} and
          \texttt{lms}; only \texttt{web-bff} is exposed here.
    \item \texttt{public}: only reverse proxy (Caddy) is exposed to the
          host network.
\end{itemize}

\textbf{Volume:}
\begin{itemize}
    \item \texttt{pgdata}: PostgreSQL data, named volume, daily backup.
    \item \texttt{minio-data}: object storage data, named volume.
    \item \texttt{qdrant-storage}: vector data, named volume.
    \item \texttt{neo4j-data}: graph data, named volume.
    \item \texttt{redis-data}: AOF persistence, named volume.
    \item \texttt{models}: BGE-M3 model weights (around 2GB), shared
          read-only mount for \texttt{document-worker} and
          \texttt{enrichment-worker}.
\end{itemize}
